\chapter{The Framework Landscape}

DCF exists within a rich ecosystem of AI agent methodologies that emerged primarily in 2024--2025 and continues to evolve. Understanding where DCF fits---and doesn't fit---clarifies its unique contribution.

\section{The Major Frameworks}

\begin{table}[H]
\centering
\small
\begin{tabular}{llll}
\toprule
\textbf{Framework} & \textbf{Focus} & \textbf{Level} & \textbf{Source} \\
\midrule
Research-Plan-Implement & Three-phase workflow & Workflow & Claude Code \\
ACE-FCA & Self-improving context & Engineering & Stanford/HumanLayer \\
12-Factor Agents & Production reliability & Engineering & Dex Horthy \\
BMAD-METHOD & Agile workflows with agents & Process & BMad Code \\
Ralph/Ralph Wiggum & Autonomous iteration loops & Operational & Geoffrey Huntley \\
Loom & Hierarchical agent architecture & Architecture & Open source \\
Awesome ecosystems & Tools, commands, integrations & Tooling & Community \\
\bottomrule
\end{tabular}
\caption{The AI Agent Methodology Landscape}
\end{table}

\section{How DCF Relates to Each}

\begin{itemize}
    \item \textbf{Research-Plan-Implement}: Symbiotic. Plan Mode provides structure; DCF provides thinking within each phase.
    \item \textbf{ACE-FCA}: Complementary. ACE automates context curation; DCF guides human evaluation of that curation.
    \item \textbf{12-Factor Agents}: Orthogonal. 12-Factor tells you how to build agents; DCF tells you how to think with agents.
    \item \textbf{BMAD-METHOD}: Overlapping. BMAD provides process structure; DCF provides thinking methodology within that process.
    \item \textbf{Ralph}: Philosophical tension. See next chapter.
    \item \textbf{Loom}: Complementary. Loom defines hierarchical structure; DCF defines how humans evaluate what that hierarchy produces.
    \item \textbf{Awesome ecosystems}: Agnostic. DCF works with any tooling.
\end{itemize}

\section{The Gap DCF Fills}

Existing frameworks address:
\begin{itemize}
    \item Research-Plan-Implement: How to STRUCTURE phased work
    \item ACE, 12-Factor: How to BUILD better agents
    \item BMAD: How to ORGANIZE agent workflows
    \item Ralph: How to AUTOMATE agent execution
    \item Loom: How to COMPOSE agent hierarchies
    \item Awesome: What TOOLS to use
\end{itemize}

DCF addresses: \textbf{How the HUMAN should THINK}---during checkpoints, reviews, approvals---the cognitive strategy layer that other frameworks leave implicit.

DCF is the \textbf{cognitive operating system} that runs on top of whatever agentic framework you choose.

\section{Emerging Validation}

The Socratic approach to human-AI collaboration is gaining independent research traction. The ``Revisited Socratic Knowledge Interaction Framework'' \cite{tonekaboni2026socratic} operationalizes the same classical elements---Elenchus, Maieutics, Aporia, and Dialectic---into a multi-agent architecture with specialized agents coordinated by a Triage Agent that routes dialogue turns based on context. Their framing of \emph{productive discomfort} (aporia as explicit design goal, not accidental byproduct) and emphasis on ``augmented intelligence'' over ``artificial intelligence'' converge with DCF's core principles.

This parallel development suggests that the Socratic method provides a natural interaction pattern for human-AI knowledge work---not because of historical accident, but because its structure addresses fundamental challenges: LLM hallucinations require human verification (elenchus), tacit knowledge must be drawn out (maieutics), and productive puzzlement motivates revision (aporia).

